\documentclass[11pt,a4paper]{article}

% ============================================================
% DSLAB submission-ready LaTeX report
% Format: 11pt, Arial-like font, one-column, max 20 pages
% ============================================================

% -----------------------------
% Font and layout
% -----------------------------
\usepackage[T1]{fontenc}
\usepackage[utf8]{inputenc}
\usepackage[margin=2.25cm]{geometry}
\usepackage{helvet}
\renewcommand{\familydefault}{\sfdefault}
\usepackage{microtype}
\usepackage{setspace}
\setstretch{1.08}
\setlength{\parindent}{0pt}
\setlength{\parskip}{0.55em}

% -----------------------------
% Tables, figures, math
% -----------------------------
\usepackage{graphicx}
\graphicspath{{./}{outputs/}}
\usepackage{booktabs}
\usepackage{tabularx}
\usepackage{array}
\usepackage{longtable}
\usepackage{makecell}
\usepackage{multirow}
\usepackage{float}
\usepackage{caption}
\usepackage{subcaption}
\usepackage[table]{xcolor}
\usepackage{enumitem}
\usepackage{amsmath,amssymb}
\usepackage{url}
\usepackage{hyperref}
\usepackage{fancyhdr}
\usepackage{titlesec}
\usepackage{tocloft}
\usepackage{listings}
\usepackage[most]{tcolorbox}

% -----------------------------
% Color palette
% -----------------------------
\definecolor{mainblue}{HTML}{1F4E79}
\definecolor{lightblue}{HTML}{EAF3FA}
\definecolor{darkgray}{HTML}{404040}
\definecolor{mediumgray}{HTML}{666666}
\definecolor{lightgray}{HTML}{F3F6F8}
\definecolor{teal}{HTML}{007A7A}
\definecolor{orange}{HTML}{F28E2B}
\definecolor{softred}{HTML}{C84C3A}
\definecolor{softgreen}{HTML}{E6F3F0}
\definecolor{softorange}{HTML}{FFF2E6}

\hypersetup{
    colorlinks=true,
    linkcolor=mainblue,
    citecolor=mainblue,
    urlcolor=mainblue
}

% -----------------------------
% Header and footer
% -----------------------------
\pagestyle{fancy}
\fancyhf{}
\fancyhead[L]{\small\textit{Life Satisfaction in OECD Countries}}
\fancyhead[R]{\small\thepage}
\renewcommand{\headrulewidth}{0.4pt}
\setlength{\headheight}{14pt}

\fancypagestyle{firstpage}{
  \fancyhf{}
  \renewcommand{\headrulewidth}{0pt}
}

% -----------------------------
% Section style
% -----------------------------
\titleformat{\section}
  {\Large\bfseries\color{mainblue}}
  {\thesection}{0.7em}{}
\titleformat{\subsection}
  {\large\bfseries\color{darkgray}}
  {\thesubsection}{0.7em}{}
\titleformat{\subsubsection}
  {\normalsize\bfseries\color{darkgray}}
  {\thesubsubsection}{0.7em}{}
\titlespacing*{\section}{0pt}{1.4em}{0.6em}
\titlespacing*{\subsection}{0pt}{1.0em}{0.3em}

% -----------------------------
% Captions, lists, code
% -----------------------------
\captionsetup{
    font=small,
    labelfont={bf,color=mainblue},
    textfont=it,
    justification=centering
}
\setlist[itemize]{leftmargin=1.2em,itemsep=0.2em,topsep=0.2em}
\setlist[enumerate]{leftmargin=1.4em,itemsep=0.2em,topsep=0.2em}
\renewcommand{\arraystretch}{1.18}

\lstset{
  basicstyle=\ttfamily\footnotesize,
  breaklines=true,
  frame=single,
  rulecolor=\color{lightblue},
  backgroundcolor=\color{lightgray},
  columns=fullflexible,
  showstringspaces=false,
  numbers=left,
  numberstyle=\tiny\color{mediumgray},
  keywordstyle=\bfseries\color{mainblue},
  commentstyle=\itshape\color{mediumgray},
  tabsize=2
}

% -----------------------------
% Custom boxes
% -----------------------------
\newtcolorbox{bluebox}{
  colback=lightblue,
  colframe=mainblue,
  arc=2mm,
  boxrule=0.7pt,
  left=8pt,
  right=8pt,
  top=6pt,
  bottom=6pt
}

\newtcolorbox{graybox}{
  colback=lightgray,
  colframe=mediumgray,
  arc=2mm,
  boxrule=0.5pt,
  left=8pt,
  right=8pt,
  top=6pt,
  bottom=6pt
}

\newcommand{\pos}[1]{\textcolor{teal}{#1}}
\newcommand{\neg}[1]{\textcolor{softred}{#1}}

\begin{document}
\thispagestyle{firstpage}

% ============================================================
% 1. Title, authors, keywords, dataset, AI declaration, abstract
% ============================================================
\begin{center}
    \vspace*{0.3cm}
    {\Large\bfseries Data Science Lab Project}\par
    \vspace{0.2cm}
    {\large February 2026 Exam Session}\par
    \vspace{0.8cm}

    \begin{tcolorbox}[colback=lightblue,colframe=mainblue,width=0.94\textwidth,arc=3mm,boxrule=0.9pt]
    \centering
    {\Huge\bfseries\color{mainblue} Life Satisfaction in OECD Countries}\par
    \vspace{0.25cm}
    {\LARGE\bfseries A Multidimensional Well-being Analysis}\par
    \vspace{0.15cm}
    {\LARGE\bfseries with a Focus on Italy}\par
    \end{tcolorbox}
\end{center}

\vspace{0.4cm}

\subsection*{Authors' Details}
\begin{table}[H]
\centering
\rowcolors{2}{lightgray}{white}
\begin{tabularx}{\textwidth}{>{\bfseries}l l l X}
\toprule
Family name & First name & Badge number & Email \\
\midrule
Dawra & Tony & 946852 & t.dawra@campus.unimib.it \\
Cavallo & Davide & 880593 & d.cavallo4@campus.unimib.it \\
Villani & Giacomo & 880819 & g.villani7@campus.unimib.it \\
\bottomrule
\end{tabularx}
\end{table}

\subsection*{Keywords}
Well-being; OECD; Life satisfaction; Social indicators; Italy

\subsection*{Name of the Dataset}
\textit{OECD Well-being Data Monitor -- Current well-being dataset}, downloaded from the OECD Data Explorer.

\subsection*{Declaration on AI Tools}
AI tools were used to help prepare reproducible R code for data cleaning, descriptive statistics, correlation analysis, visualization, construction of dimension scores, and simple regression modelling. The final data selection, execution of the analysis, verification of numerical outputs, and interpretation of the results, formulate the question of the research question, organization of the report, and the drafting of the explanatory text remain the responsibility of the authors. AI tools were not used to fabricate data, numerical results, or references.

\subsection*{Abstract}
\begin{bluebox}
This study investigates which dimensions of well-being are most associated with life satisfaction across OECD countries, with a specific focus on Italy compared with the OECD average excluding Italy. The analysis uses the OECD Well-being Data Monitor, Current well-being dataset, covering 2010--2025, 38 OECD countries, and 21 selected measures. Since indicators are not always available for the same year, the latest available observation is retained for each country and measure. Life satisfaction is the outcome variable, while explanatory indicators are grouped into seven dimensions: economic security, work quality, housing, health, social connections, civic engagement, and environment and safety. The methodology combines descriptive statistics, missing-value analysis, correlations, standardized dimension scores, Italy-OECD comparison, and a simple regression model. Results show that social connections have the strongest dimension-level association with life satisfaction ($r=0.453$), followed by environment and safety ($r=0.305$) and work quality ($r=0.300$). Italy reports lower life satisfaction than the OECD benchmark and performs particularly below it in social connections, civic engagement, work quality, and environment and safety.
\end{bluebox}

\newpage
\tableofcontents
\newpage

% ============================================================
\section{Research Question}
% ============================================================

\subsection{Statement}

\begin{bluebox}
\textbf{Research question.} Which dimensions of well-being are most associated with life satisfaction across OECD countries, and how does Italy compare with the OECD average excluding Italy?
\end{bluebox}

The dependent variable is \textit{life satisfaction}, measured on a 0--10 scale. The explanatory variables are not treated as isolated predictors only; instead, they are grouped into conceptual dimensions of well-being. This makes the analysis more interpretable and more consistent with the purpose of the project: starting from a relevant social question and using data analysis to answer it in a transparent and realistic way.

The explanatory dimensions are economic security, work quality, housing, health, social connections, civic engagement, and environment and safety. The aim is not only to check whether richer countries are more satisfied, but also to investigate whether broader dimensions of daily life are associated with subjective well-being.

\subsection{Motivation and Stakeholders}

The OECD argues that measuring progress requires looking beyond GDP and considering broader economic, social, and environmental outcomes for people \cite{oecd_beyond_gdp}. This is the main motivation of the study: income and production are important, but they do not fully describe whether people live well.

The question is relevant for several stakeholders. Policy makers and governments may use well-being evidence to identify whether priorities should focus more on income, employment, health, housing, social relations, safety, or civic participation. Students and young workers are especially affected by employment quality, housing affordability, and social support. Researchers can use the results as a descriptive comparison of OECD countries. Italy is used as a country focus because it provides a concrete interpretation of the OECD benchmark.

This project is exploratory. It does not claim to prove that one variable causes higher or lower life satisfaction. The objective is to describe the strongest associations and discuss them carefully, especially given the limitations of aggregate country-level data.

% ============================================================
\section{The Data}
% ============================================================

\subsection{Data Description}

The dataset used in this project is the \textit{OECD Well-being Data Monitor -- Current well-being dataset}. The Current well-being data explorer allows users to explore measures of current well-being by country and compare them across OECD members \cite{oecd_current}. The OECD has 38 member countries, which is also the number of countries selected in this project \cite{oecd_members}.

The original downloaded file contained 8,594 observations and 30 columns. After cleaning, filtering, and selecting the latest available observation for each country and measure, the final analytical dataset contained 38 OECD countries and 21 numeric well-being indicators plus the country name. Therefore, the final wide dataset had 38 rows and 22 columns.

\begin{table}[H]
\centering
\caption{Dataset filters used in the OECD Data Explorer.}
\rowcolors{2}{lightgray}{white}
\begin{tabular}{>{\bfseries}ll}
\toprule
Filter & Selection \\
\midrule
Time period & 2010--2025 \\
Reference area & 38 OECD countries \\
Measure & 21 well-being measures \\
Age & Total \\
Sex & Total \\
Education level & Total \\
Observation used & Latest available value \\
\bottomrule
\end{tabular}
\end{table}

The variable \textit{life satisfaction} is the outcome variable. The remaining variables are grouped into seven dimensions of well-being. The grouping is used to make the interpretation clearer and to avoid reducing the analysis to a list of isolated indicators.

\begin{table}[H]
\centering
\caption{Selected measures and conceptual well-being dimensions.}
\rowcolors{2}{lightgray}{white}
\begin{tabularx}{\textwidth}{>{\bfseries}p{0.23\textwidth} X}
\toprule
Dimension & Selected measures \\
\midrule
Outcome & Life satisfaction \\
Economic security & Households and NPISHs net adjusted disposable income per capita; median net wealth; household disposable income below the relative income poverty line \\
Work quality & Employment rate; labour market insecurity; long hours in paid work; job strain \\
Housing & Housing affordability; housing cost overburden; households living in overcrowded conditions \\
Health & Life expectancy at birth; perceived health as positive \\
Social connections & Social support; lack of social support \\
Civic engagement & Having a say in government; voter turnout \\
Environment and safety & Exposed to air pollution; exposure to extreme temperature; feeling safe at night; road deaths \\
\bottomrule
\end{tabularx}
\end{table}

\subsection{Strengths and Weaknesses of the Data}

The main strength of the dataset is that it comes from the OECD and is designed for cross-country comparison. It covers several dimensions of well-being, making it appropriate for a multidimensional analysis rather than a purely economic one. The OECD's \textit{How's Life? 2024} report also emphasizes the use of many indicators to assess current well-being outcomes, inequalities, and resources for future well-being \cite{oecd_howslife}.

There are also limitations. First, not all variables are available for all countries. Second, different indicators may refer to different latest years. Third, the analysis is based on country-level aggregates and cannot explain individual differences within countries. Finally, the associations observed in the data should not be interpreted as causal effects.

\begin{figure}[H]
\centering
\includegraphics[width=0.78\textwidth]{plot_missing_values.png}
\caption{Missing values for the selected indicators. The plot highlights which variables have the weakest data availability.}
\label{fig:missing}
\end{figure}

% ============================================================
\section{Methodology}
% ============================================================

\subsection{Overall Logic Thread of the Study}

The methodology follows five steps. First, the data are cleaned and reshaped from long format to wide format. Second, descriptive statistics and missing values are examined. Third, each selected measure is correlated with life satisfaction. Fourth, the measures are grouped into well-being dimensions and converted into standardized dimension scores. Fifth, Italy is compared with the OECD average excluding Italy, and a simple regression model is estimated using the strongest dimensions.

Because the dataset covers 2010--2025 and not every indicator is available for every country in every year, using a single reference year would have produced many missing values. For this reason, the analysis selects the latest available observation for each country and each measure. This maximizes country coverage while keeping the analysis as recent as possible.

Therefore, the analysis should be interpreted as a latest-available cross-sectional comparison, not as a comparison based on one single common year for all indicators.
\begin{lstlisting}[language=R,caption={Selection of latest available observation.}]
df_latest_long <- df_selected %>%
  group_by(country, measure_key) %>%
  arrange(desc(year), .by_group = TRUE) %>%
  slice(1) %>%
  ungroup()
\end{lstlisting}

\subsection{Tools: Which and Why}

The analysis was conducted in RStudio. R was selected because it is appropriate for reproducible data cleaning, descriptive statistics, visualization, correlation analysis, and regression modelling. The main packages used were \texttt{tidyverse}, \texttt{janitor}, \texttt{skimr}, \texttt{broom}, and \texttt{ggplot2}.

No global imputation was applied. Descriptive statistics and correlations were computed using available observations. The regression model was estimated on complete cases for the selected variables. This choice avoids creating artificial values in a dataset with only 38 countries and keeps the analysis transparent.

\subsubsection{Correlation Analysis}

For each numeric indicator $X_j$, the Pearson correlation with life satisfaction $Y$ is computed as:

\[
 r_{Y,X_j} = \frac{\sum_i (Y_i - \bar{Y})(X_{ij} - \bar{X}_j)}{\sqrt{\sum_i (Y_i - \bar{Y})^2}\sqrt{\sum_i (X_{ij} - \bar{X}_j)^2}}.
\]

Positive values indicate that higher values of an indicator are associated with higher life satisfaction, while negative values indicate the opposite.

\subsubsection{Dimension Scores}

Since the variables are expressed in different units, each measure was standardized using a z-score:

\[
 z_{ij} = \frac{x_{ij} - \bar{x}_j}{s_j}.
\]

For indicators where higher values represent worse conditions, such as poverty, job strain, lack of social support, air pollution, road deaths, and housing cost overburden, the standardized value was multiplied by $-1$. Therefore, higher dimension scores always represent better conditions. The score for each dimension was then computed as the average of the available standardized components.

\subsubsection{Regression Model}

A simple regression model was used as an interpretative tool. The three dimensions with the strongest absolute correlations with life satisfaction were included: social connections, environment and safety, and work quality. The estimated model is:

\[
 LS_i = \beta_0 + \beta_1 SC_i + \beta_2 ES_i + \beta_3 WQ_i + \epsilon_i.
\]

Here, $SC_i$ denotes the social-connections score, $ES_i$ the environment-and-safety score, and $WQ_i$ the work-quality score. The model is not intended to prove causality. It is used to check whether the strongest dimensions remain informative when considered together.

% ============================================================
\section{Results}
% ============================================================

\subsection{Steps of the Study: Description and Essential Numbers}

Table~\ref{tab:desc} presents descriptive statistics for selected indicators. Across the OECD countries with available data, average life satisfaction is 7.33 on a 0--10 scale, with a standard deviation of 0.54. The observed values range from 5.70 to 8.45. This shows that most OECD countries are concentrated within a relatively narrow range of life satisfaction, although relevant differences between countries still exist.

\begin{table}[H]
\centering
\caption{Descriptive statistics for selected indicators.}
\label{tab:desc}
\rowcolors{2}{lightgray}{white}
\begin{tabular}{lrrrrr}
\toprule
Indicator & n & Mean & SD & Min & Max \\
\midrule
Life satisfaction & 34 & 7.33 & 0.54 & 5.70 & 8.45 \\
Employment rate & 38 & 78.72 & 5.65 & 59.24 & 87.61 \\
Social support & 38 & 89.83 & 3.85 & 79.82 & 97.66 \\
Life expectancy & 38 & 81.13 & 2.50 & 75.50 & 84.30 \\
Housing cost overburden & 36 & 18.26 & 10.13 & 3.76 & 48.53 \\
Relative income poverty & 37 & 11.48 & 3.90 & 5.52 & 21.02 \\
Job strain & 33 & 27.59 & 6.87 & 13.83 & 47.94 \\
Perceived health positive & 37 & 68.62 & 10.96 & 37.10 & 85.60 \\
Road deaths & 36 & 4.61 & 2.21 & 1.45 & 12.24 \\
\bottomrule
\end{tabular}
\end{table}

The first analytical step was to compute correlations between life satisfaction and the selected individual measures. Table~\ref{tab:correlations} reports the strongest individual associations.

\begin{table}[H]
\centering
\caption{Strongest individual correlations with life satisfaction.}
\label{tab:correlations}
\rowcolors{2}{lightgray}{white}
\begin{tabular}{llr}
\toprule
Measure & Dimension & Correlation \\
\midrule
job strain & Work quality & {-0.541} \\
perceived positive health & Health & \pos{0.516} \\
road deaths & Environment and safety & {-0.494} \\
social support & Social connections & \pos{0.461} \\
relative income poverty & Economic security & {-0.457} \\
lack of social support & Social connections & {-0.442} \\
labour-market insecurity & Work quality & {-0.350} \\
employment rate & Work quality & \pos{0.297} \\
air pollution exposure & Environment and safety & {-0.289} \\
extreme temperature exposure & Environment and safety & {-0.184} \\
say in government & Civic engagement & \pos{0.181} \\
disposable income & Economic security & \pos{0.089} \\
\bottomrule
\end{tabular}
\end{table}

\begin{figure}[H]
\centering
\includegraphics[width=0.86\textwidth]{plot_measure_correlations.png}
\caption{Strongest individual correlations with life satisfaction. Teal bars indicate positive associations; red-orange bars indicate negative associations.}
\label{fig:measurecorr}
\end{figure}

At the dimension level, social connections show the strongest association with life satisfaction, with a correlation of $0.453$. Environment and safety have the second strongest association, with a correlation of $0.305$, followed by work quality with $0.300$. Economic security also shows a positive relationship, but weaker than social connections, with a correlation of $0.245$.

\begin{table}[H]
\centering
\caption{Dimension-level correlations with life satisfaction.}
\label{tab:dimcorr}
\rowcolors{2}{lightgray}{white}
\begin{tabular}{lr}
\toprule
Dimension & Correlation \\
\midrule
Social connections & \pos{0.453} \\
Environment and safety & \pos{0.305} \\
Work quality & \pos{0.300} \\
Economic security & \pos{0.245} \\
Health & \pos{0.162} \\
Housing & {-0.069} \\
Civic engagement & {-0.008} \\
\bottomrule
\end{tabular}
\end{table}

\begin{figure}[H]
\centering
\includegraphics[width=0.82\textwidth]{plot_dimension_correlations.png}
\caption{Correlation between standardized well-being dimension scores and life satisfaction.}
\label{fig:dimcorr}
\end{figure}

Job strain is the strongest individual correlate of life satisfaction. Figure~\ref{fig:jobstrain} shows the country-level relationship between job strain and life satisfaction. The fitted line slopes downward, confirming the negative association.

\begin{figure}[H]
\centering
\includegraphics[width=0.84\textwidth]{plot_top_measure_scatter.png}
\caption{Scatterplot of life satisfaction and job strain. The fitted regression line indicates a negative association.}
\label{fig:jobstrain}
\end{figure}

Italy is compared with the OECD average excluding Italy. This choice avoids including Italy in the benchmark against which it is evaluated. Since dimension scores are standardized, values above zero indicate better-than-average conditions relative to the OECD distribution, while values below zero indicate weaker conditions.

\begin{table}[H]
\centering
\caption{Italy compared with the OECD average excluding Italy.}
\label{tab:italy}
\rowcolors{2}{lightgray}{white}
\begin{tabular}{lrrr}
\toprule
Indicator & Italy & OECD excl. Italy & Difference \\
\midrule
Life satisfaction & 7.100 & 7.334 & {-0.234} \\
Health & 0.784 & -0.052 & \pos{0.835} \\
Economic security & 0.026 & -0.059 & \pos{0.086} \\
Housing & -0.158 & -0.005 & {-0.154} \\
Environment and safety & -0.295 & 0.000 & {-0.295} \\
Work quality & -0.517 & -0.019 & {-0.498} \\
Civic engagement & -0.657 & 0.018 & {-0.675} \\
Social connections & -0.935 & 0.025 & {-0.960} \\
\bottomrule
\end{tabular}
\end{table}

\begin{figure}[H]
\centering
\includegraphics[width=0.84\textwidth]{plot_italy_vs_oecd_dimensions_excluding_italy.png}
\caption{Italy compared with the OECD average excluding Italy across standardized well-being dimensions. Positive values indicate better-than-average performance relative to the OECD distribution, while negative values indicate weaker performance. The blue bars represent the OECD average excluding Italy.}
\label{fig:italy}
\end{figure}

Finally, a regression model was estimated using the three strongest dimensions from the correlation analysis. The model uses 34 complete observations because four countries are missing life satisfaction values or some variables used in the model.

\begin{table}[H]
\centering
\caption{Regression coefficients. Dependent variable: life satisfaction.}
\label{tab:reg}
\rowcolors{2}{lightgray}{white}
\begin{tabular}{lrrrr}
\toprule
Term & Estimate & SE & t & p-value \\
\midrule
Intercept & 7.321 & 0.089 & 82.286 & 0.000 \\
Social connections & 0.229 & 0.114 & 2.005 & 0.054 \\
Environment and safety & 0.059 & 0.253 & 0.232 & 0.818 \\
Work quality & -0.035 & 0.208 & -0.166 & 0.869 \\
\bottomrule
\end{tabular}
\end{table}

\begin{table}[H]
\centering
\caption{Regression model summary.}
\label{tab:regsummary}
\rowcolors{2}{lightgray}{white}
\begin{tabular}{lr}
\toprule
Statistic & Value \\
\midrule
Observations & 34 \\
$R^2$ & 0.207 \\
Adjusted $R^2$ & 0.127 \\
Residual standard error & 0.501 \\
F-statistic & 2.606 \\
Model p-value & 0.070 \\
\bottomrule
\end{tabular}
\end{table}

\subsection{Comments and Interpretation}

The results answer the research question in a clear but cautious way. Among the well-being dimensions, social connections are most strongly associated with life satisfaction. This finding is important because it shows that subjective well-being is not simply a function of income or wealth. Economic security is positively related to life satisfaction, but its dimension-level correlation is weaker than that of social connections.

The second group of relevant dimensions is environment and safety and work quality. These dimensions are also positively associated with life satisfaction when bad indicators are reversed. At the individual-measure level, job strain is particularly important, which suggests that working conditions can influence well-being beyond the simple employment rate.

Housing and civic engagement show weak associations in this dataset. This does not imply that housing and civic engagement are unimportant in general. Instead, it suggests that the selected indicators, the latest-observation design, and the small cross-country sample do not reveal a strong linear relationship with life satisfaction for those dimensions.

The Italy comparison strengthens the interpretation. Italy is relatively strong in health, close to the benchmark in economic security, but weaker in social connections, civic engagement, work quality, environment and safety, and housing. Since social connections are the strongest dimension-level correlate of life satisfaction, Italy's low social-connection score is especially relevant.

The regression explains about 20.7\% of the cross-country variation in life satisfaction. The coefficient for social connections is positive and close to conventional statistical significance ($p=0.054$). The coefficients for environment and safety and work quality are not statistically significant in the joint model. Therefore, the regression should be interpreted as supporting evidence, not as a definitive explanatory model.

The small negative coefficient for work quality should not be interpreted as evidence that better work quality reduces life satisfaction. In the simple dimension-level correlation, work quality is positively associated with life satisfaction, while the sign changes only in the joint regression, likely because the dimensions overlap and the sample contains only 34 complete observations.
% ============================================================
\section{Conclusion}
% ============================================================

This project analysed which dimensions of well-being are most associated with life satisfaction across OECD countries and compared Italy with the OECD average excluding Italy. The analysis used the OECD Well-being Data Monitor, selecting 21 measures for 38 countries over the period 2010--2025 and retaining the latest available observation for each country and measure.

The main result is that social connections have the strongest dimension-level association with life satisfaction. Environment and safety and work quality follow, while economic security is positively associated but not dominant. At the individual-measure level, job strain, perceived positive health, road deaths, social support, poverty, and lack of social support are among the strongest correlates.

For Italy, the analysis shows a lower life satisfaction score than the OECD average excluding Italy. Italy's main strength is health, while its main weaknesses are social connections, civic engagement, work quality, and environment and safety. These results suggest that policies aimed at improving quality of life should not be limited to income growth, but should also address social support, job strain, safety, and broader living conditions.

The main strength of the study is its use of internationally comparable OECD data and an interpretable methodology. The main weakness is that the analysis is based on aggregate country-level data, which prevents individual-level conclusions and causal claims. Future research could improve the analysis by using individual-level survey data, adding a real time-series component, or focusing on specific groups such as young people, low-income households, or workers.

% ============================================================
\section{References}
% ============================================================

\begin{thebibliography}{9}

\bibitem{oecd_monitor}
OECD. \textit{OECD Well-being Data Monitor}. Available at: \url{https://www.oecd.org/en/data/tools/well-being-data-monitor.html}.

\bibitem{oecd_beyond_gdp}
OECD. \textit{Well-being and beyond GDP}. Available at: \url{https://www.oecd.org/en/topics/policy-issues/well-being-and-beyond-gdp.html}

\bibitem{oecd_current}
OECD. \textit{Current well-being Data Explorer}. Available at: \url{https://data-explorer.oecd.org/vis?fs[0]=Topic,1%7CSociety%23SOC%23%7CWell-being%20and%20beyond%20GDP%23SOC_WEL%23&pg=0&fc=Topic&bp=true&snb=26&df[ds]=dsDisseminateFinalDMZ&df[id]=DSD_CWB%40DF_CWB&df[ag]=OECD.WISE.CWB&df[vs]=1.0&dq=..A1&lom=LASTNPERIODS&lo=5&to[TIME_PERIOD]=false}.

\bibitem{oecd_members}
OECD. \textit{Members and partners}. Available at: \url{https://www.oecd.org/en/about/members-partners.html}.

\bibitem{oecd_howslife}
OECD. \textit{How's Life? 2024: Well-being and Resilience in Times of Crisis}. OECD Publishing, Paris, 2024. Available at: \url{https://www.oecd.org/en/publications/2024/11/how-s-life-2024_bdcf2f9f.html}.



\end{thebibliography}

% ============================================================
\appendix
\section{Appendix}
% ============================================================

\subsection{Variables Description}

\begin{longtable}{p{0.28\textwidth}p{0.22\textwidth}p{0.42\textwidth}}
\caption{Complete list of selected variables and dimensions.}\label{tab:appendixvars}\\
\toprule
Variable & Dimension & Interpretation \\
\midrule
\endfirsthead
\toprule
Variable & Dimension & Interpretation \\
\midrule
\endhead
life satisfaction & Outcome & Average evaluation of life on a 0--10 scale \\
income & Economic security & Disposable income per capita \\
wealth & Economic security & Median net wealth \\
poverty & Economic security & Share below relative income poverty line; higher is worse \\
employment rate & Work quality & Share of working-age population employed \\
labour-market insecurity & Work quality & Expected earnings loss from unemployment; higher is worse \\
long hours & Work quality & Share working long paid hours; higher is worse \\
job strain & Work quality & Share experiencing high job strain; higher is worse \\
housing affordability & Housing & Housing affordability indicator; higher is better \\
housing cost overburden & Housing & Housing cost overburden; higher is worse \\
overcrowding & Housing & Share living in overcrowded conditions; higher is worse \\
life expectancy & Health & Life expectancy at birth \\
perceived health positive & Health & Share reporting positive health \\
social support & Social connections & Share reporting social support \\
lack social support & Social connections & Share lacking social support; higher is worse \\
say in government & Civic engagement & Share reporting a say in government \\
voter turnout & Civic engagement & Voter turnout \\
air pollution & Environment and safety & Share exposed to air pollution; higher is worse \\
extreme temperature & Environment and safety & Exposure to extreme temperature; higher is worse \\
feeling safe & Environment and safety & Share feeling safe at night \\
road deaths & Environment and safety & Road deaths; higher is worse \\
\bottomrule
\end{longtable}

\subsection{Full Descriptive Statistics}

\begin{table}[H]
\centering
\small
\caption{Descriptive statistics for all numeric indicators.}
\rowcolors{2}{lightgray}{white}
\begin{tabular}{lrrrrr}
\toprule
Variable & n & Mean & SD & Min & Max \\
\midrule
air pollution & 37 & 94.95 & 11.71 & 47.85 & 100.00 \\
employment rate & 38 & 78.72 & 5.65 & 59.24 & 87.61 \\
extreme temperature & 38 & 6.52 & 13.23 & 0.00 & 44.19 \\
feeling safe & 38 & 72.99 & 12.47 & 37.40 & 90.73 \\
housing affordability & 37 & 80.14 & 2.86 & 73.76 & 85.63 \\
housing cost overburden & 36 & 18.26 & 10.13 & 3.76 & 48.53 \\
income & 34 & 39811.26 & 9576.36 & 19128.00 & 63190.00 \\
job strain & 33 & 27.59 & 6.87 & 13.83 & 47.94 \\
labour-market insecurity & 35 & 5.11 & 4.26 & 0.99 & 21.68 \\
lack social support & 38 & 9.63 & 3.88 & 2.17 & 19.92 \\
life expectancy & 38 & 81.13 & 2.50 & 75.50 & 84.30 \\
life satisfaction & 34 & 7.33 & 0.54 & 5.70 & 8.45 \\
long hours & 36 & 6.47 & 5.98 & 0.49 & 25.91 \\
perceived health positive & 37 & 68.62 & 10.96 & 37.10 & 85.60 \\
poverty & 37 & 11.48 & 3.90 & 5.52 & 21.02 \\
road deaths & 36 & 4.61 & 2.21 & 1.45 & 12.24 \\
say in government & 32 & 29.43 & 10.68 & 12.10 & 58.17 \\
social support & 38 & 89.83 & 3.85 & 79.82 & 97.66 \\
voter turnout & 38 & 70.95 & 11.02 & 46.64 & 90.67 \\
wealth & 30 & 217036.87 & 141907.85 & 61894.00 & 853921.00 \\
overcrowding & 35 & 11.53 & 8.64 & 0.59 & 31.57 \\
\bottomrule
\end{tabular}
\end{table}

\subsection{Reproducibility Code Excerpt}

\begin{lstlisting}[language=R,caption={Core R workflow used for the analysis.}]
library(tidyverse)
library(janitor)
library(broom)

# Import data
df_raw <- read_csv("OECD_WELLBEING_DATA.csv", show_col_types = FALSE) %>%
  clean_names()

# Latest observation for each country and measure
df_latest_long <- df_selected %>%
  group_by(country, measure_key) %>%
  arrange(desc(year), .by_group = TRUE) %>%
  slice(1) %>%
  ungroup()

# Wide dataset for correlations and regression
df_wide <- df_latest_long %>%
  select(country, measure_key, value) %>%
  pivot_wider(names_from = measure_key, values_from = value)

# Italy vs OECD benchmark excluding Italy
oecd_average_row <- analysis_scores %>%
  filter(str_to_lower(country) != "italy") %>%
  summarise(across(where(is.numeric), ~ mean(., na.rm = TRUE)))

# Regression with the three strongest dimensions
model <- lm(life_satisfaction ~ dim_social_connections +
              dim_environment_safety + dim_work_quality,
            data = analysis_scores)

tidy(model)
glance(model)
\end{lstlisting}

\end{document}
